\documentclass[10pt]{article}
\usepackage[margin=1in]{geometry}
\usepackage{amsmath,amssymb,amsthm}
\usepackage[T1]{fontenc}
\usepackage{lmodern}
\usepackage[colorlinks=true,linkcolor=blue,citecolor=blue,urlcolor=blue]{hyperref}
\newtheorem{theorem}{Theorem}
\newtheorem{lemma}{Lemma}
\newcommand{\R}{\mathbb R}
\title{A telescoping proof of the focal antipedal sum invariant}
\author{Alec Kriebel\\Independent researcher\\
  \small\href{https://orcid.org/0009-0001-9320-500X}{ORCID 0009-0001-9320-500X}}
\date{6 October 2026\quad---\quad Version 1.0}
\begin{document}
\maketitle
\begin{abstract}
Reznik, Garcia and Koiller observed that the sums of the distances from
the two foci of an ellipse to their respective antipedal vertices agree
along every closed elliptic billiard orbit. We prove this equality for
every regular closed orbit tangent to a strictly nested confocal ellipse,
including odd primitive periods, star orbits and repeated traversals.
The difference of the two ordinary positive antipedal distances on each
edge is a fixed coefficient times the vertical displacement of that edge.
Closure therefore gives the equality without a parity restriction.
\end{abstract}

\section{Statement, construction and attribution}
Let
\[
 E:\ \frac{x^2}{a^2}+\frac{y^2}{b^2}=1,\qquad
 E_\lambda:\ \frac{x^2}{a^2-\lambda}
                  +\frac{y^2}{b^2-\lambda}=1,
 \qquad a>b>0,\quad 0<\lambda<b^2.
\]
Put $c=\sqrt{a^2-b^2}$ and $f_\sigma=(\sigma c,0)$ for
$\sigma\in\{+1,-1\}$. A regular closed billiard/Poncelet orbit here is a
cyclic list $P_1,\ldots,P_N$ on $E$, with distinct consecutive endpoints,
whose edges are tangent to $E_\lambda$ and whose continuation at each
vertex uses the other of its two caustic tangents. Thus immediate
retracing is excluded. The list may self-intersect, have any admissible
winding, or repeat a shorter orbit. Set $P_{N+1}=P_1$.

For a focus $f_\sigma$, define $Q_{\sigma,i}$ as the intersection of
the full lines through $P_i$ and $P_{i+1}$ perpendicular to
$P_i-f_\sigma$ and $P_{i+1}-f_\sigma$, respectively. These are the
vertices of the antipedal of the original orbit polygon about that focus.
We use ordinary Euclidean distances
$q_{\sigma,i}=\lVert Q_{\sigma,i}-f_\sigma\rVert$.

\begin{theorem}\label{thm:main}
Every regular closed orbit just defined has finite, strictly positive
antipedal distances and satisfies
\begin{equation}\label{eq:sum}
 \sum_{i=1}^N q_{+,i}=\sum_{i=1}^N q_{-,i}.
\end{equation}
In particular the ratio of these two sums is $1$ for every member of
every admitted Poncelet family.
\end{theorem}

This is the observation labelled $k_{603}$ by Reznik--Garcia--Koiller
in \cite[\S\S3.5,3.7 and Table 7]{RGKpre} and
\cite[\S3.7 and Table 7, p.~349]{RGKjournal}. Their table records its
experimental discovery in May 2020 and marks its proof as unknown there.
That dated marker is not a claim about the present literature. The
observation and antipedal construction are theirs, and the underlying
negative-pedal/polarity geometry is classical \cite[Arts.~121--122]{Salmon}.
The focal sideline-height product used below is also known
\cite[Lemma 4.2]{BT}. Central inversion already proves the equality for
centrally paired orbits; ordinary telescoping is a familiar method.
The contribution proposed here is the identification of the ordinary
norm difference with a coordinate increment and the resulting full-period
proof, in particular for odd primitive periods.

Neither individual sum is asserted constant as an orbit varies. The
claim concerns the original polygon, not the outer tangent polygon,
pedal feet, side perimeters or signed radial lengths. Hyperbolic caustics,
degenerate caustics and two-bounce limits are outside its hypotheses.

\section{Positive antipedal distance and orientation}
\begin{lemma}\label{lem:distance}
For noncollinear $A,B,f$, the antipedal intersection $Q$ is unique and
\begin{equation}\label{eq:distance}
 \lVert Q-f\rVert=\frac{\lVert A-f\rVert\lVert B-f\rVert}{h},
 \qquad h=\operatorname{dist}(f,\operatorname{line}(A,B))>0.
\end{equation}
\end{lemma}
\begin{proof}
Translate $f$ to zero and write $A=(r,0)$,
$B=(s\cos\alpha,s\sin\alpha)$ with $r,s>0$ and
$\sin\alpha\ne0$. The defining lines give
$Q\cdot A=r^2$, $Q\cdot B=s^2$, hence
\[
 \lVert Q\rVert^2
 =\frac{r^2+s^2-2rs\cos\alpha}{\sin^2\alpha}
 =\frac{\lVert A-B\rVert^2}{\sin^2\alpha}.
\]
The triangle area identity gives
$h=rs|\sin\alpha|/\lVert A-B\rVert$. Taking the positive square root
proves \eqref{eq:distance}.
\end{proof}

The whole orbit can be oriented with the caustic on the left of each
directed edge. Indeed, a tangent line to the strictly inner ellipse meets
$E$ in two points, and its contact lies strictly between them. At a
vertex the rays toward the two contact points bound a cone containing
the caustic, so the caustic lies on opposite sides of those two directed
rays. Reversing the incoming ray and continuing on the other tangent
therefore preserves its side. The nonretracing rule propagates one
consistent side throughout the orbit; reverse the list if it is the
right side. This is a local argument and does not require a convex
polygon or any ordering of its vertices around $E$.

For such a directed edge $A B$, take its outward unit normal
$n=(u,v)$ and tangent $t=(-v,u)$. Its equation is
$n\cdot X=\rho$, with $u^2+v^2=1$ and $\rho>0$; the caustic and
origin lie on the side $n\cdot X<\rho$. Put
\begin{equation}\label{eq:K}
 K=a^2u^2+b^2v^2,\qquad \rho^2=K-\lambda.
\end{equation}
The second equality is the support-function condition for tangency
to $E_\lambda$. The chord endpoints are
\begin{equation}\label{eq:chord}
 A=M-\ell t,\quad B=M+\ell t,\qquad
 M=\left(\frac{a^2\rho u}{K},\frac{b^2\rho v}{K}\right),
 \quad \ell=\frac{ab\sqrt\lambda}{K}>0.
\end{equation}
To check this, substitute $M+z t$ in the outer ellipse equation:
the linear term vanishes, the constant is $\rho^2/K$, and the
quadratic coefficient is $K/(a^2b^2)$. Thus $z=\pm\ell$.

The two focal heights are
\begin{equation}\label{eq:heights}
 h_\sigma=\rho-\sigma cu,
 \qquad h_+h_-=\rho^2-c^2u^2=b^2-\lambda>0.
\end{equation}
Their sum is $2\rho>0$, so both heights are strictly positive.
They are therefore the ordinary perpendicular distances of the foci
to the line, not signed substitutes. Equivalently, both foci lie strictly
inside $E_\lambda$, since $c^2<a^2-\lambda$.
No edge is focal, and Lemma~\ref{lem:distance} applies to both foci.

\section{The edge identity and closure}
For a point $(x,y)$ on $E$, its distance to $f_\sigma$ is
$a-\sigma cx/a$: squaring verifies the identity, and this expression
is at least $a-c>0$. From \eqref{eq:chord},
\[
 x_A+x_B=\frac{2a^2\rho u}{K},\qquad
 x_Ax_B=\frac{a^2(\rho^2-b^2v^2)}{K}.
\]
Multiplying the two positive endpoint distances gives
\begin{equation}\label{eq:product}
 R_\sigma:=\lVert A-f_\sigma\rVert\lVert B-f_\sigma\rVert
 =\frac{a^2h_\sigma^2+b^2\lambda}{K}.
\end{equation}
For completeness, the product before simplification is
$a^2-\sigma c(x_A+x_B)+(c^2/a^2)x_Ax_B$; inserting the preceding
two expressions and $\rho^2=K-\lambda$ yields \eqref{eq:product}.
Lemma~\ref{lem:distance} now gives positive distances
\begin{equation}\label{eq:q}
 q_\sigma=\frac{1}{K}
     \left(a^2h_\sigma+\frac{b^2\lambda}{h_\sigma}\right).
\end{equation}
Since $h_+-h_-=-2cu$ and
$h_+^{-1}-h_-^{-1}=2cu/(b^2-\lambda)$, subtraction gives
\begin{equation}\label{eq:difference}
 q_+-q_-=\frac{2cu}{K}
      \left(-a^2+\frac{b^2\lambda}{b^2-\lambda}\right).
\end{equation}
But \eqref{eq:chord} also gives
$y_B-y_A=2ab\sqrt\lambda\,u/K$. Hence
\begin{equation}\label{eq:coboundary}
 \boxed{q_+-q_-=\Gamma(y_B-y_A)},\qquad
 \Gamma=\frac{c}{ab\sqrt\lambda}
       \left(-a^2+\frac{b^2\lambda}{b^2-\lambda}\right).
\end{equation}
All denominators are positive. No division by $u$ or by the vertical
displacement occurs, so horizontal chords are included. The coefficient
vanishes at $\lambda=a^2b^2/(a^2+b^2)$ and changes sign there;
neither fact affects positivity of $q_\sigma$.
Formula \eqref{eq:coboundary} uses the chosen caustic-on-left orientation.
Reversing endpoints leaves $q_+-q_-$ unchanged but negates $y_B-y_A$;
in the reverse orientation its displacement coefficient is $-\Gamma$.

\begin{proof}[Proof of Theorem~\ref{thm:main}]
Use the consistent left orientation above. The same $\lambda$ and
$\Gamma$ apply to every edge. Summing \eqref{eq:coboundary} gives
\[
 \sum_i(q_{+,i}-q_{-,i})
 =\Gamma\sum_i(y_{i+1}-y_i)=0
\]
by closure. Equations \eqref{eq:heights} and \eqref{eq:distance} show
that every summand is positive and finite, so the denominator of the
ratio is nonzero. No parity condition is used. Repeated traversal simply
repeats the summands. Reversing an edge leaves its two antipedal
intersections unchanged, hence reversing the orbit preserves the equality.
\end{proof}

\section{Verification and status}
The accompanying standard-library Python verifier independently solves
the two antipedal line equations on rational chords. It checks the norm,
positive focal heights, confocal tangency and the edge identity exactly;
it also covers horizontal chords, zero and opposite signs of $\Gamma$,
reversal and a closed repeated four-cycle. Deliberately incorrect
coefficient, height, norm sign and manuscript-hash controls must fail.
All guards remain active under Python optimization. Its recorded counts
are in the support files. These finite computations supplement the
written proof for all real parameters and every admitted period.

A bounded primary-source and mechanism audit completed on 6 October 2026
located no earlier full covering proof. In particular, focal inverse
perimeters \cite{RGR} concern different summands, and the complete
journal final \cite{GKR}, compared with its precursor, provides spatial
averages for other observables rather than this finite antipedal norm
difference. This is not an absolute-first or continued-open-status
guarantee. The support audit summary retains the explicit final-book
versus author-source boundary, unread historical and experimental
materials, and finite search/version limits.

\paragraph{Disclosure.}
AI tools were used extensively in solving, deriving, drafting, literature
checking, code generation, verification and adversarial checking.
The independent checks described in the support package include AI
agents and are not conventional human peer review. This manuscript is
unrefereed and has not received conventional human peer review.
The manuscript and original support files are licensed under
\href{https://creativecommons.org/licenses/by/4.0/}{CC BY 4.0}; cited
third-party works retain their own terms and are not redistributed.

\begin{thebibliography}{9}
\bibitem{RGKpre} D. Reznik, R. Garcia and J. Koiller,
\emph{Eighty New Invariants of N-Periodics in the Elliptic Billiard},
arXiv:2004.12497v11 (29 October 2020),
\href{https://arxiv.org/abs/2004.12497v11}{arXiv:2004.12497v11}.
\bibitem{RGKjournal} D. Reznik, R. Garcia and J. Koiller,
\emph{Fifty New Invariants of N-Periodics in the Elliptic Billiard},
Arnold Mathematical Journal 7 (2021), 341--355,
\href{https://doi.org/10.1007/s40598-021-00174-y}{doi:10.1007/s40598-021-00174-y}.
\bibitem{Salmon} G. Salmon,
\emph{A Treatise on the Higher Plane Curves}, 3rd ed.,
Hodges, Foster, and Figgis, Dublin, 1879, Arts.~121--122, pp.~105--107.
\bibitem{BT} M. Bialy and S. Tabachnikov,
\emph{Dan Reznik's identities and more},
European Journal of Mathematics 8 (2022), 1341--1354,
\href{https://doi.org/10.1007/s40879-020-00428-7}{doi:10.1007/s40879-020-00428-7}.
\bibitem{RGR} P. Roitman, R. Garcia and D. Reznik,
\emph{New Invariants of Poncelet--Jacobi Bicentric Polygons},
Arnold Mathematical Journal 7 (2021), 619--637,
\href{https://doi.org/10.1007/s40598-021-00188-6}{doi:10.1007/s40598-021-00188-6}.
\bibitem{GKR} R. Garcia, J. Koiller and D. Reznik,
\emph{Estimating Elliptic Billiard Invariants with Spatial Integrals},
Journal of Dynamical and Control Systems 29 (2023), 757--767
(online 10 August 2022),
\href{https://doi.org/10.1007/s10883-022-09608-y}{doi:10.1007/s10883-022-09608-y}.
\end{thebibliography}
\end{document}
